\section*{问题重述}
本题要求在复杂场景下形成从多模态观测到情感预测、再到可核验证据的分析链。问题一针对附件1的原始媒体，分别提取文本、语音和视觉特征，将不同采样粒度的观测映射到公共时间轴，并保留每个有效位置的来源。问题二使用附件2已提供的对齐数值特征，预测三类情感及连续情感强度，重点考察局部连续时间段缺失时的性能。问题三在确定的预测模型上分析模态贡献、模态交互和关键特征区间，并对附件4的无标签样本给出预测与解释。

三个问题在时间位置、有效掩码和证据可追溯性方面相互衔接，但各自使用的官方附件不同。问题一构建的附件1特征并未替代问题二的附件2输入；问题三的附件4推理也不提供有标签性能评价。本文据此分别报告特征构建的完整性、附件2验证集上的预测与缺失场景表现，以及附件4上的解释结果。
